% Select--compress--reinvest: controlled visual-token allocation study.
% arXiv / ARR working copy of main.tex. Do not edit main.tex from this pass.
% arXiv: \usepackage[preprint]{acl} (authors visible).
% ARR:   \usepackage[review]{acl}  (anonymous; also drop the GitHub URL).
\documentclass[11pt]{article}
\usepackage[preprint]{acl}
\usepackage{times}
\usepackage{latexsym}
\usepackage{booktabs}
\usepackage{graphicx}
\usepackage{amsmath}
\usepackage{amssymb}
\usepackage{xcolor}
\usepackage{multirow}
\usepackage{url}

\graphicspath{{../figures/}}

\title{A Controlled Study of Keyframe Selection and Visual-Token Allocation in Video MLLMs}

\author{Prakhar Khatri \\
  Independent Researcher \\
  \texttt{prakharkhatri123@gmail.com}}

\begin{document}
\maketitle

\begin{abstract}
Long-video multimodal large language models can encode only a small fraction of a video, so a fixed visual-token budget is spent first on which frames to keep and then on how finely each frame is encoded.
Published comparisons often change the scorer, selection rule, resolution policy, and answerer together, which makes those two decisions hard to isolate.
We freeze a LongCLIP question-stem scorer and a matched evaluation harness, and compare six training-free selectors at eight keyframes on LongVideoBench, Video-MME, and LVBench.
Query-aware selection is the largest lever.
Classical Orthogonal Matching Pursuit (OMP), used here as a cheap reference rather than a proposed method, improves over uniform sampling by 5.7, 5.9, and 11.8 percentage points on the three benchmarks; LDDR's stage-1 selector remains within one point, with no detectable paired difference on Video-MME or LVBench.
OMP with eight frames also outperforms sixteen-frame uniform sampling on hour-long videos while using half as many frames.
Holding timestamps fixed, reducing the per-frame spatial budget by roughly half changes benchmark-level accuracy by at most 0.44 points; on pooled LongVideoBench long videos, the compressed arm is equivalent in accuracy to full-resolution encoding within a $\pm3$-point margin.
Reinvesting those tokens in sixteen lower-resolution keyframes then improves over eight full-resolution frames by 2.36 points on LongVideoBench's long bins and 3.04 points on LVBench.
Thus, select relevant keyframes first; spatial compression is useful mainly as a way to buy broader temporal coverage at fixed cost.
\end{abstract}

\section{Introduction}

Long-video question answering begins with keyframe sampling.
A video MLLM cannot usually encode every frame, so its input pipeline must identify a small set of moments that preserves the evidence needed for the question.
Uniform temporal sampling remains a common default, while recent systems replace it with query-aware retrieval, temporal coverage, diversity, or learned evidence scores~\citep{tang2025aks,sun2025mdp3,zhu2025focus,wang2026evidential,peng2026qca}.
Other systems jointly adapt frame selection and spatial resolution~\citep{zhang2025qframe,chen2026lddr,wang2026dafs} or prune tokens after frames enter the vision stack~\citep{vistallm2026,moprune2026}.
Together, these methods show that which frames enter the answerer is a first-order design decision.

Their end-to-end improvements leave a narrower keyframe-sampling question unresolved: under a fixed scorer, frame budget, and answerer, which query-aware selection rule provides the most useful evidence?
Published comparisons do not directly answer this question because the scorer, prompt, frame count, resolution policy, and answerer often change together.
A selector can appear better because it uses a stronger encoder; a dynamic-resolution system can improve because it adds frames rather than because it allocates resolution well.
After isolating the selection rule, a second question follows: can spatial compression extend the selected keyframe set without increasing visual-token cost?

We answer both questions with matched counterfactuals built around a single LongCLIP question-stem scorer and a shared evaluation harness.
The study proceeds in three steps.
First, we compare uniform sampling, top-$k$, AKS, a controlled replay of FOCUS, OMP, and LDDR's stage-1 MinMax selector at $k{=}8$.
Second, we freeze the timestamps and reduce only the spatial budget per frame.
Third, we reinvest the saved tokens by moving from eight full-resolution frames to sixteen compressed frames at equal or lower measured visual-token cost.
Classical OMP~\citep{pati1993omp} is a strong inexpensive selector in this controlled comparison; the contribution is the empirical allocation result that its use exposes, rather than the pursuit algorithm itself.

The evidence places keyframe selection at the center of the resulting select--compress--reinvest strategy.
OMP improves substantially over uniform sampling on long videos and can outperform uniform baselines that consume twice as many frames.
Spatial compression is a supporting intervention: at fixed timestamps, halving the per-frame budget barely changes the observed accuracy, with formal $\pm3$-point equivalence established on pooled LongVideoBench long videos.
Using that released budget for additional selected timestamps yields a further two-to-three-point gain in the best-supported settings.
The benefit is not universal across answerers and benchmarks, and later picks trade coverage against off-topic novelty.

\paragraph{Contributions.}
\begin{itemize}
  \item We provide a matched keyframe-selection comparison of six training-free rules under one scorer, prompt boundary, frame budget, and answerer harness. OMP emerges as a strong query-focused selector and remains within one point of LDDR's stage-1 rule.
  \item We show that selective sampling can dominate larger uniform frame sets: OMP uses half as many frames while significantly outperforming uniform sampling in two long-video cross-budget comparisons.
  \item We show how to extend a selected keyframe set under fixed cost. Roughly halving per-frame spatial budget preserves the observed accuracy, while reinvesting the saved tokens in additional keyframes improves long-video QA in the strongest supported settings.
  \item Paired tests, two answerer families, three benchmarks, residual diagnostics, and a qualitative failure audit delimit when query-focused diversity supplies useful coverage and when it drifts toward irrelevant content.
\end{itemize}

\section{Related Work}

\paragraph{Query-aware frame selection.}
Training-free selectors typically combine three objectives: query relevance, coverage of the timeline, and diversity within the selected set.
AKS recursively partitions the timeline to preserve coverage~\citep{tang2025aks}; FOCUS treats temporal clips as arms in a budgeted bandit~\citep{zhu2025focus}; and MDP3 models relevance, list-wise diversity, and sequentiality with a DPP-based dynamic program~\citep{sun2025mdp3}.
AdaRD-Key and Adaptive Greedy likewise treat relevance--diversity as a training-free objective~\citep{adardkey2025,adaptivegreedy2026}.
Recent segment-based methods make this structure explicit.
QCA assigns segment quotas using relevance and content deviation~\citep{peng2026qca}, while EFS forms event-like segments and refines query-relevant anchors with adaptive maximal marginal relevance~\citep{chen2026efs}.
ReQuest adds uncertainty-triggered computation and adaptive temporal suppression~\citep{kim2026request}; query-conditioned evidential sampling instead learns a frame-level estimate of conditional evidence~\citep{wang2026evidential}.
These approaches differ not only in subset selection but also in scorer, candidate pool, training requirements, and frame budget.

\paragraph{Selection with spatial adaptation.}
Q-Frame couples CLIP-based query-aware selection with multi-resolution scaling so that more frames fit within a compute limit~\citep{zhang2025qframe}.
LDDR combines a linear-DPP selector with Group-DPP importance for frame retention and dynamic resolution~\citep{chen2026lddr}.
DAFS extracts relevance from MLLM attention and jointly optimizes candidate-pool size and per-frame token budget~\citep{wang2026dafs}.
These systems establish the value of joint allocation, but their coupled designs do not isolate whether a gain comes from the selected timestamps, the spatial policy, or the additional temporal coverage.
Our study complements them by intervening on these axes separately.

\paragraph{Visual-token reduction.}
Token-level methods reduce cost after or during visual encoding.
AdaCodec learns retain, compress, and drop decisions~\citep{adacodec2026}; AdaptToken assigns group budgets from answerer uncertainty~\citep{adapttoken2026}; Vista-LLM performs query-guided pruning before inference~\citep{vistallm2026}; and MoPrune uses scene structure and motion to retain informative tokens~\citep{moprune2026}.
This work asks a coarser but experimentally separable question: before token-level pruning, does a long-video answerer benefit more from sharper retained frames or from seeing additional moments?
A concurrent allocator, AdaAlloc~\citep{an2026adalloc}, splits a fixed budget between a low-resolution overview and high-resolution local evidence.
We do not claim to outperform such allocators; a reconstruction of LDDR's stage-2 importance already beats our residual-proportional rule.

\section{Controlled Keyframe Selection and Budget Interventions}
\label{sec:method}

\begin{figure*}[t]
\centering
\includegraphics[width=\textwidth]{fig_protocol.pdf}
\caption{Controlled protocol.
(a) A 1\,fps candidate pool is encoded once with LongCLIP using only the question stem; every $k{=}8$ selector consumes the same cached embeddings and scores.
(b) Selected timestamps are then frozen while spatial budget is reduced.
(c) The saved visual tokens are reinvested by comparing eight full-resolution frames with sixteen compressed frames at equal or lower measured token cost.
FOCUS$^\star$ denotes a replay of the published selection schedule on dense LongCLIP scores, not its original budgeted ITM scorer.}
\label{fig:protocol}
\end{figure*}

\subsection{Controlled interventions}

The experimental unit is a benchmark question paired across two input policies.
We study three interventions in sequence:
\textbf{selection}, which changes timestamps while holding $k$, scorer, and resolution fixed;
\textbf{compression}, which holds timestamps fixed and changes the per-frame spatial budget; and
\textbf{reinvestment}, which spends the saved spatial budget on additional timestamps.
This design measures the marginal value of each decision rather than comparing full systems whose components differ simultaneously.

\subsection{Candidate pool and shared scorer}

Videos are decoded at 1\,fps.
Each candidate frame and the question stem are encoded with LongCLIP~\citep{zhang2024longclip}; answer options never enter selection.
This boundary matters empirically: using the fused question-and-options string changed roughly 40\% of selected frames on LongVideoBench-600\,s and could reverse the score ordering.
Frame embeddings and stem similarities are cached once, and every selector consumes the same cache.

\subsection{Selection rules}

The matched comparison contains uniform sampling, cosine top-$k$, AKS~\citep{tang2025aks}, FOCUS$^\star$~\citep{zhu2025focus}, OMP, and LDDR-select~\citep{chen2026lddr}.
FOCUS$^\star$ replays the published clip-bandit schedule on the dense LongCLIP vector; Appendix~\ref{app:focus} states the resulting fidelity boundary.
LDDR-select contains only the MinMax stage-1 selector, not the full dynamic-resolution pipeline.

Let $E\in\mathbb{R}^{N\times d}$ contain L2-normalized frame embeddings and let $q_0$ be the normalized stem embedding.
OMP greedily selects a frame correlated with the current residual and then projects the query away from the selected span:
\begin{equation}
\begin{aligned}
b_r &= \arg\max_{i\notin B_{r-1}} e_i^\top q_{r-1},\\
q_r &= q_0-\operatorname{Proj}_{\operatorname{span}\{e_j:j\in B_r\}}(q_0),
\end{aligned}
\label{eq:omp}
\end{equation}
where $B_r=B_{r-1}\cup\{b_r\}$.
In this embedding space, the projection mainly suppresses directions already represented by earlier picks; Section~\ref{sec:mechanism} tests rather than assumes a sparse-reconstruction interpretation.

\subsection{Spatial budget and reinvestment}

For the fixed-timestamp intervention, frames are resized before the model processor, whose internal resize is disabled for Qwen3-VL.
The main compressed arm uses a residual-proportional resize schedule with a target mean spatial fraction of about 0.53 (D@53).
We also compare a flat split at the same mean fraction and a reconstruction of LDDR's stage-2 Group-DPP importance on identical OMP-8 timestamps.
The flat control tests whether OMP rank alone is a useful resolution priority; the LDDR reconstruction tests whether a different allocation shape can matter.

The reinvestment arm selects $k{=}16$ timestamps and applies an approximately 50\% per-frame cap, distinct from the 0.53 mean used when timestamps stay fixed.
An empirical token audit reproduces Qwen's smart-resize rules from each video's dimensions.
The sixteen-frame arm consumes 0.996 and 0.984 times the tokens of the eight-frame full-resolution arm in the 600\,s and 3600\,s LongVideoBench bins, respectively.
Thus, the LongVideoBench reinvestment comparison is conservative with respect to total visual tokens.
Video-MME and LVBench use the same resize rule; we do not report a separate token audit for those benches.

\section{Evaluation Setup}
\label{sec:setup}

We evaluate LongVideoBench~\citep{wu2024longvideobench} (official validation set, $n{=}1337$), Video-MME~\citep{fu2025videomme} ($n{=}2700$), and LVBench~\citep{lvbench2025} ($n{=}1549$).
LongVideoBench is additionally stratified into 15, 60, 600, and 3600\,s duration bins.
The primary answerer is Qwen3-VL-8B-Instruct~\citep{qwen3vl2025}, evaluated through lmms-eval~\citep{zhang2024lmmseval} with greedy decoding, temperature zero, and subtitles disabled.
InternVL3-2B and InternVL3-8B~\citep{zhu2025internvl3} provide a second open-model family.
GPT-5-mini checks selected contrasts with the same saved timestamps.

\begin{table}[t]
\centering
\small
\begin{tabular}{@{}lp{0.58\columnwidth}@{}}
\toprule
Benchmarks & LongVideoBench (1337), Video-MME (2700), LVBench (1549) \\
Primary answerer & Qwen3-VL-8B, greedy decoding, subtitles off \\
Transfer checks & InternVL3 2B/8B; GPT-5-mini \\
Shared scorer & LongCLIP, question stem only, 1\,fps pool \\
Frame budgets & $k{=}8$ for selection; $k{=}16$ for reinvestment \\
\bottomrule
\end{tabular}
\caption{Evaluation grid. Main claims concern paired accuracy differences under matched inputs, not absolute leaderboard position.}
\label{tab:setup}
\end{table}

Accuracy comparisons use exact two-sided McNemar tests~\citep{mcnemar1947} on paired correctness outcomes.
Because failure to reject a difference does not establish equivalence, the compression analysis also reports two one-sided tests (TOST) and 90\% confidence intervals~\citep{schuirmann1987}.
The $\pm2$, $\pm3$, and $\pm4$ percentage-point margins were examined after the runs rather than preregistered; the text therefore reports the narrowest supported margin and the adjacent unsupported margin.
Tests are contrast-specific and uncorrected for multiplicity.

\section{Results}
\label{sec:results}

\subsection{Query-aware keyframe selection is the largest lever}
\label{sec:selection}

Table~\ref{tab:t1} compares selectors at $k{=}8$ under the shared scorer.
OMP improves over uniform sampling by 5.69 points on LongVideoBench, 5.85 on Video-MME, and 11.81 on LVBench.
The corresponding paired evidence is strongest in the LongVideoBench 3600\,s bin ($p{=}5.9\times10^{-4}$), Video-MME ($p{=}5.1\times10^{-10}$), and LVBench ($p{=}9.5\times10^{-17}$).
LDDR-select remains within one point of OMP across all three benchmarks; paired tests detect no difference on Video-MME or LVBench ($p{=}.70$ and $.57$).

\begin{table*}[t]
\centering
\small
\begin{tabular}{lccc}
\toprule
Method & LongVideoBench & Video-MME & LVBench \\
\midrule
Uniform & .5654 ($-$5.69) & .5637 ($-$5.85) & .3454 ($-$11.81) \\
Top-$k$ & .6028 ($-$1.95) & .5704 ($-$5.18) & .4319 ($-$3.16) \\
\textbf{OMP} & .6223 & \textbf{.6222} & .4635 \\
AKS & .6021 ($-$2.02) & .5785 ($-$4.37) & .4280 ($-$3.55) \\
FOCUS$^\star$ & .5819 ($-$4.04) & .5578 ($-$6.44) & .3983 ($-$6.52) \\
LDDR-select & \textbf{.6320} (${+}0.97$) & .6193 ($-$0.29) & \textbf{.4693} (${+}0.58$) \\
\bottomrule
\end{tabular}
\caption{Matched selector comparison at $k{=}8$ with Qwen3-VL-8B. Cells are accuracy; parentheses give the difference from OMP in percentage points. FOCUS$^\star$ is the controlled schedule replay described in Appendix~\ref{app:focus}; LDDR-select is stage~1 only.}
\label{tab:t1}
\end{table*}

\begin{figure}[t]
\centering
\includegraphics[width=\columnwidth]{fig2_long_video_effect.pdf}
\caption{OMP's gain over uniform sampling grows with video duration on LongVideoBench. At 15\,s, uniform, top-$k$, and OMP have identical accuracy (.7249); the significant gains occur at 600 and 3600\,s.}
\label{fig:longvideo}
\end{figure}

The duration split locates the effect.
At 15\,s, uniform, top-$k$, and OMP are identical; OMP exceeds uniform by 7.8 points at 600\,s ($p{=}.0022$) and 7.5 points at 3600\,s ($p{=}5.9\times10^{-4}$).
Selection also compensates for frame scarcity across budgets.
On 3600\,s videos, OMP with eight frames reaches .5461 versus .4770 for sixteen uniform frames, a 6.9-point gain while using half as many frames ($p{=}.0011$).
On 600\,s videos, OMP-16 reaches .6578 versus .6044 for uniform-32, a 5.3-point gain at half the frame count ($p{=}.018$).
These contrasts support selective sparse sampling more directly than a same-$k$ comparison alone.

The strength of simpler relevance rules is benchmark-dependent.
Top-$k$, AKS, and FOCUS$^\star$ do not separate detectably from uniform on Video-MME, whereas OMP and LDDR-select do.
On LongVideoBench, OMP's advantage over top-$k$ is smaller and often not detectable within individual bins.
The controlled conclusion is therefore specific: combining query relevance with suppression of already represented directions is consistently strong in this harness, while the value of any particular simpler heuristic does not transfer automatically.

\subsection{Full spatial detail is often unnecessary}
\label{sec:compression}

Table~\ref{tab:t2} freezes timestamps and changes only spatial budget.
Across the three Qwen benchmark aggregates, moving to the compressed arm changes accuracy by at most 0.44 points.
The uniform-timestamp control on pooled LongVideoBench long videos also changes by only 0.31 points ($p{=}.784$), showing that the stability is not confined to OMP-selected frames.

\begin{table}[t]
\centering
\small
\resizebox{\columnwidth}{!}{%
\begin{tabular}{lccc}
\toprule
Answerer / benchmark & Full & Compressed & $\Delta$ (pp) \\
\midrule
Qwen / LongVideoBench & .6223 & .6253 & ${+}0.30$ \\
Qwen / Video-MME & .6222 & .6178 & $-0.44$ \\
Qwen / LVBench & .4635 & .4674 & ${+}0.39$ \\
GPT-5-mini / LVBench & .4900 & .4926 & ${+}0.26$ ($p{=}.84$) \\
Uniform / LVB long pool & .5061 & .5092 & ${+}0.31$ ($p{=}.784$) \\
\bottomrule
\end{tabular}}
\caption{Fixed-timestamp spatial compression. The Qwen rows use the saved approximately half-budget arm; the GPT row uses a pixel-ratio control and is not a Qwen token count. The uniform control uses an approximately 46\% realized per-frame budget.}
\label{tab:t2}
\end{table}

Only the pooled LongVideoBench long-bin comparison has a formal equivalence result for Qwen.
For $n{=}976$, the compressed OMP arm is 0.72 points higher with a 90\% interval of $[-0.60,+2.03]$.
This interval lies inside $\pm3$ points (TOST $p{=}.0022$) but narrowly misses $\pm2$ points ($p{=}.054$).
InternVL3-8B provides a stronger spatial intervention: reducing 24 tiles to 8 changes pooled accuracy by $-0.10$ points, with a 90\% interval of $[-1.66,+1.45]$, establishing $\pm2$-point equivalence ($p{=}.022$).
The small LVBench and Video-MME point differences in Table~\ref{tab:t2} are descriptive stability checks, not formal equivalence claims.

\begin{figure}[t]
\centering
\includegraphics[width=\columnwidth]{fig_claim_a.pdf}
\caption{LongVideoBench 3600\,s accuracy on fixed OMP-8 timestamps across spatial budgets. The flat 50\% control and the residual-proportional schedule are indistinguishable, so OMP rank alone does not provide a useful resolution policy.}
\label{fig:claim-a}
\end{figure}

The allocation shape still matters when it contains a stronger importance signal.
Our residual-proportional schedule does not outperform a flat split at the same mean budget.
In contrast, a reconstruction of LDDR's stage-2 Group-DPP importance improves over a flat split by 1.84 points on the pooled long bins ($p{=}.0198$, 95\% CI $[+0.37,+3.32]$).
That contrast is from one same-environment reconstruction rather than the authors' released implementation, and its 3600\,s component is not individually significant.
The supported conclusion is consequently narrower than ``resolution allocation is irrelevant'': coarse compression is inexpensive, a simple OMP-rank allocation adds no value, and a more structured importance rule can add value.

\subsection{Reinvesting pixels in frames improves accuracy}
\label{sec:reinvest}

The reinvestment intervention compares eight full-resolution frames with sixteen compressed frames at equal or lower measured token cost.
On LongVideoBench, the overall gain is 2.24 points; pooling the 600 and 3600\,s bins gives 2.36 points ($p{=}.0346$), with nearly identical positive effects in both bins.
LVBench improves by 3.04 points ($p{=}.0009$), Video-MME by 1.56 points, and GPT-5-mini on LVBench by 2.39 points ($p{=}.039$).

\begin{table}[t]
\centering
\small
\resizebox{\columnwidth}{!}{%
\begin{tabular}{lccc}
\toprule
Answerer / benchmark & $k{=}8$ full & $k{=}16$ compressed & $\Delta$ (pp) \\
\midrule
Qwen / LongVideoBench & .6223 & .6447 & ${+}2.24$ \\
Qwen / LVB long pool & .5820 & .6055 & ${+}2.36$ ($p{=}.0346$) \\
Qwen / Video-MME & .6222 & .6378 & ${+}1.56$ \\
Qwen / LVBench & .4635 & .4939 & ${+}3.04$ ($p{=}.0009$) \\
GPT-5-mini / LVBench & .4900 & .5139 & ${+}2.39$ ($p{=}.039$) \\
\bottomrule
\end{tabular}}
\caption{Reinvestment at equal or lower visual-token cost on the audited LongVideoBench bins. The compressed arm approximately doubles temporal coverage; Video-MME and LVBench use the same resize rule without a separate token audit.}
\label{tab:t3}
\end{table}

\begin{figure}[t]
\centering
\includegraphics[width=0.80\columnwidth]{fig_claim_b.pdf}
\caption{LongVideoBench 600\,s and 3600\,s pooled: sixteen compressed OMP frames outperform eight full-resolution OMP frames at a measured token ratio no greater than one.}
\label{fig:claim-b}
\end{figure}

The data do not establish that reinvestment requires OMP.
Under uniform sampling, the same eight-to-sixteen intervention has a positive 1.33-point estimate but is not significant ($p{=}.283$).
The paired selector-by-budget interaction is 1.02 points ($p{=}.4435$): OMP gains on 109 items for which uniform does not, while uniform gains on 97 items for which OMP does not.
This is evidence against a large OMP-specific interaction, not proof of identical effects.
The safe interpretation is that temporal reinvestment works with OMP and is directionally consistent under uniform sampling; selector specificity remains unresolved at approximately the one-point scale.

\subsection{What later OMP picks contribute}
\label{sec:mechanism}

OMP's residual trace explains why additional frames can help after the query-directed signal has saturated.
On both LongVideoBench long bins, the residual norm is 0.972 of its initial value after the first pick and 0.967 after pick eight; it does not decrease further by pick sixteen.
Residual correlation with the selected frame falls from .233 at pick one to .004 at pick eight and approximately zero at pick sixteen.
Yet the later selected frames retain an original-query cosine near .20.
They therefore add little new query direction while increasing the chance that the answerer observes another relevant moment.

\begin{figure*}[t]
\centering
\includegraphics[width=\textwidth]{fig1_residual_inert.pdf}
\caption{LongCLIP geometry on LongVideoBench-600\,s ($n{=}411$).
(a) Residual correlation vanishes while residual norm changes little; the line connects logged picks 1, 8, and 16 and is not a dense trace.
(b) Within-video frames occupy a coherent cone, while the text query is nearly orthogonal to the frame dictionary. In this space, OMP mainly suppresses directions already represented by earlier selections.}
\label{fig:residual}
\end{figure*}

More coverage is useful only while it remains relevant.
An exploratory visual audit inspected 93 selected OMP failures at $k{=}8$ across the 600 and 3600\,s bins.
The dominant observed failure was off-topic novelty: 20 of 41 inspected 600\,s cases and 35 of 52 inspected 3600\,s cases spent several picks on visually distinctive but question-irrelevant segments, title cards, or transitions.
Conversely, sequence-of-scenes and cross-scene tracking questions sometimes benefited from OMP's spread because dense top-$k$ sampling omitted one of several required scenes.
The direct $k{=}8$ versus $k{=}16$ category results are mixed, so the audit supports a mechanism hypothesis rather than a universal scene-versus-temporal rule.
The design implication is a relevance-constrained diversity objective: preserve spread within plausible evidence regions while preventing late picks from drifting to unrelated chapters.

\subsection{The same timestamps transfer unevenly}
\label{sec:transfer}

Table~\ref{tab:internvl} reuses the same OMP and uniform timestamps with InternVL3.
The selection gain replicates on LongVideoBench and LVBench at both model sizes.
Video-MME is different: the 2B model benefits at every duration, whereas the 8B model benefits only on short videos and is flat on medium and long videos despite receiving different frames.
The same 8B model gains 10.14 points on LVBench, ruling out a model-size-only explanation.

\begin{table}[t]
\centering
\small
\resizebox{\columnwidth}{!}{%
\begin{tabular}{lrrr}
\toprule
Benchmark & $n$ & IVL3-2B $\Delta$ & IVL3-8B $\Delta$ \\
\midrule
Video-MME short & 900 & ${+}6.00$ ($p{=}1.2\times10^{-4}$) & ${+}5.44$ ($p{=}1.9\times10^{-4}$) \\
Video-MME medium & 900 & ${+}3.33$ ($p{=}.046$) & ${+}0.33$ ($p{=}.891$) \\
Video-MME long & 900 & ${+}4.22$ ($p{=}.0062$) & ${+}0.11$ ($p{=}1.000$) \\
LVBench & 1549 & ${+}9.04$ ($p{=}7.2\times10^{-11}$) & ${+}10.14$ ($p{=}1.2\times10^{-13}$) \\
LongVideoBench & 1337 & ${+}4.71$ ($p{=}1.6\times10^{-4}$) & ${+}6.51$ ($p{<}10^{-4}$) \\
\bottomrule
\end{tabular}}
\caption{OMP minus uniform accuracy on identical timestamps and tile budgets with InternVL3. Deltas are percentage points.}
\label{tab:internvl}
\end{table}

The allocation results also transfer.
On pooled LongVideoBench long videos, InternVL3-8B preserves accuracy after a threefold tile reduction ($-0.10$ points, $p{=}1.000$) and improves by 3.38 points when the same total tile budget is redistributed from 8 to 24 uniformly sampled frames ($p{=}.0099$).
The latter effect is concentrated in the 600\,s bin, which reinforces the broader conclusion: selection and allocation effects are properties of an answerer--benchmark pair, not constants of a timestamp set.

\section{Discussion}
\label{sec:discussion}

The experiments identify keyframe selection as the primary decision under a fixed visual-token budget.
A small query-focused set outperforms uniform sampling, sometimes even when the uniform input contains twice as many frames.
Compression changes few predictions in the tested Qwen and InternVL settings, and reinvesting the released budget in temporal coverage produces a smaller additional gain.
Selection should therefore precede spatial allocation.

This decomposition sharpens the interpretation of joint allocation systems such as Q-Frame, LDDR, and DAFS~\citep{zhang2025qframe,chen2026lddr,wang2026dafs}.
Their end-to-end gains combine several causal routes.
Our controls find no value from a simple priority-proportional resolution rule, a positive effect from reconstructed LDDR-style importance, and the clearest downstream gain from temporal reinvestment.
Future methods should report separate timestamp, resolution, and frame-count interventions alongside the joint result.

OMP's relevance-plus-decorrelation behavior is a strong operating point under LongCLIP, although its residual does not literally reconstruct the text query.
Late picks can extend multi-scene coverage after the text residual flattens, but unconstrained diversity also rewards irrelevant novelty.
Relevance floors, event boundaries, and stronger cross-modal scorers are natural next steps.

Selection remains answerer- and benchmark-dependent: InternVL3-8B is nearly insensitive to OMP on medium and long Video-MME but highly sensitive on LVBench.
Claims about long-video keyframe selection should therefore span multiple duration regimes and more than one answerer family.

\section{Limitations}
\label{sec:limitations}

\paragraph{Controlled components are narrower than full systems.}
FOCUS$^\star$ replays the published selection schedule on dense LongCLIP scores rather than reproducing the original budgeted ITM scorer.
LDDR-select contains stage~1 only, and the stage-2 allocation contrast is reconstructed from the paper rather than the authors' code.

\paragraph{Equivalence is established in one primary setting.}
The formal Qwen equivalence result applies to pooled LongVideoBench 600 and 3600\,s examples within a post-hoc $\pm3$-point margin; it misses $\pm2$ points.
Video-MME compression lacks a paired test, and the small LVBench deltas are descriptive rather than equivalence evidence.
The selector-by-budget interaction is underpowered for effects near one point.

\paragraph{Mechanism evidence is encoder- and sample-specific.}
Residual geometry and Gram statistics use frozen LongCLIP stem embeddings, primarily on LongVideoBench-600\,s.
They may not transfer to SigLIP, BLIP-ITM, MLLM-attention scorers, or subtitle-aware selection.
The model-assisted failure audit examines selected failures rather than a random sample, so its counts are exploratory.

\paragraph{Evaluation scope is incomplete.}
Subtitles are disabled, making subtitle-anchored questions partly inaccessible to every visual selector.
Some controls ran on a second GPU and software stack; only within-environment contrasts are interpreted when drift was observed.
GPT-5-mini compression uses a pixel ratio rather than Qwen token accounting, and some GPT selector results survive only as aggregates.

\section{Conclusion}
\label{sec:conclusion}

Query-aware keyframe selection is the dominant lever in this controlled long-video study.
Under matched controls, OMP substantially outperforms uniform sampling, remains close to LDDR's stage-1 selector, and can beat uniform inputs with twice as many frames.
Roughly half of the tested per-frame spatial budget can then be removed with little observed change, with $\pm3$-point equivalence on pooled LongVideoBench long videos.
Reinvesting those tokens in additional keyframes yields a further two-to-three-point gain in the strongest settings.
Select relevant keyframes first, compress them conservatively, and widen coverage without abandoning relevance.

\section{Ethical Considerations}
\label{sec:ethics}

The experiments use existing public video-QA benchmarks and frozen open or API models.
No new human subjects, personal data, videos, or annotations were collected.
Question stems, but not answer options, enter the selector.
Benchmark videos remain governed by their original licenses and are not redistributed with the planned artifacts.

\section{Data and Code Availability}

Evaluation code, selector implementations, analysis scripts, per-item predictions permitted by the benchmark licenses, and aggregate result tables are available at \url{https://github.com/codeprakhar25/omp-keyframe-sampling}.
The release excludes benchmark video files and API credentials.

\section{Author Contributions}

Prakhar Khatri: Conceptualization, Methodology, Software, Validation, Formal Analysis, Investigation, Data Curation, Visualization, Writing--Original Draft, and Writing--Review \& Editing.

\section{Funding}

This research received no external funding.

\section{Competing Interests}

The author declares no competing interests.

\section{AI Assistance Disclosure}

Generative AI tools assisted code development, manuscript editing, and an exploratory visual failure audit.
The author directed the analyses, verified the reported numerical results against saved artifacts, adjudicated the scientific claims, and takes responsibility for the final manuscript.

\bibliography{references}

\appendix

\section{FOCUS$^\star$ schedule replay}
\label{app:focus}

Every row in Table~\ref{tab:t1}, including FOCUS$^\star$, consumes the same dense 1\,fps LongCLIP stem-score vector.
The original FOCUS method instead treats scoring as a budgeted online process in which clip-arm pulls purchase ITM evidence~\citep{zhu2025focus}.
Our replay preserves the temporal clip schedule and default hyperparameters but reveals already-computed LongCLIP scores to the bandit.
It therefore tests the schedule under a matched scorer, not the efficiency or accuracy of the full FOCUS pipeline.

The replay uses 16\,s clips, three coarse pulls per clip, $\alpha{=}0.25$, a pull budget equal to half of the candidate pool, and approximately $k/4$ surviving clips.
Unobserved rewards are interpolated from the nearest observed frame, and the random seed is fixed per item.
This boundary is why the row carries a star and why its result is not used to make a claim about FOCUS under its native ITM protocol.

\section{Question categories do not define a stable selector regime}
\label{app:qtype}

LongVideoBench provides question-category tags.
We compare OMP with uniform sampling on temporally referred categories (the \texttt{T*} group) and all remaining categories without relabeling the data.

\begin{table}[h]
\centering
\small
\resizebox{\columnwidth}{!}{%
\begin{tabular}{llrccc}
\toprule
Bin & Slice & $n$ & Uniform & OMP & $\Delta$ (McNemar) \\
\midrule
600\,s & temporal \texttt{T*} & 148 & .5608 & .5811 & ${+}2.03$, $p{=}.74$ \\
600\,s & other & 264 & .5492 & .6591 & ${+}10.98$, $p{=}5.2\times10^{-4}$ \\
3600\,s & temporal \texttt{T*} & 222 & .4369 & .5270 & ${+}9.01$, $p{=}.013$ \\
3600\,s & other & 342 & .4942 & .5585 & ${+}6.43$, $p{=}.021$ \\
\bottomrule
\end{tabular}}
\caption{OMP minus uniform on official LongVideoBench categories. Percentage-point gains differ by duration, and the category interaction is not significant in either long bin.}
\label{tab:qtype}
\end{table}

The temporal slice appears resistant to selection at 600\,s but improves at 3600\,s.
Many \texttt{T*} stems refer to subtitles or speech while the selector is vision-only and subtitles are disabled, providing one plausible reason for instability.
Because the interaction is not significant and category winners change across bins, we do not infer a question-type router from these results.

\section{Selector variants}
\label{app:negatives}

\begin{figure}[t]
\centering
\includegraphics[width=\columnwidth]{fig4_negative_sweep.pdf}
\caption{Coverage-gated selector variants against OMP at $k{=}8$ on the 600 and 3600\,s LongVideoBench bins. The shaded $\pm1.5$-point band is a visual aid, not a statistical criterion.}
\label{fig:negatives}
\end{figure}

The variant sweep changes only the subset rule on the same LongCLIP cache.
MMR underperforms; query-blind DPP ($\beta{=}0$) collapses toward uniform sampling; and query-weighted DPP, residual floors, and partial orthogonalization remain in a band around OMP.
No tested DPP variant significantly exceeds OMP in either long bin.
These nulls do not imply that all diversity objectives are equivalent.
They show that, in this scorer geometry and budget range, changing the selection rule within the tested relevance--diversity family yields less movement than replacing uniform sampling or reallocating pixels to frames.

\section{Exploratory failure audit}
\label{app:failure}

Two independent model-assisted visual passes inspected OMP failures at $k{=}8$ through a frame viewer: 41 cases spanning 16 categories in the 600\,s bin and 52 cases spanning 17 categories in the 3600\,s bin.
Cases were selected to cover categories and both OMP-only and shared failures; the sample is not random.
Reviewers compared selected timestamps, rendered frames, the question, answer, and predictions.
Primary tags included missed evidence, temporal offset, answerer error despite adequate frames, and questions not answerable from visual frames.

Both passes identified the same dominant OMP-specific pattern: visually distinctive but irrelevant content consumed several of eight picks.
Examples included unrelated chapters in compilations, different storylines in long narrative video, title or subscribe cards, and dark transitions.
The opposite failure appeared on sequence-of-scenes and cross-scene tracking questions, where dense top-$k$ sampling covered one relevant moment but omitted another required scene.
The audit therefore motivates relevance-constrained diversity, but it cannot estimate population failure rates or prove that adding frames benefits a particular category.

\end{document}
